% ============================================================================
%  A crystallographic restriction for multiplier cyclic sieving on necklaces
%
%  Compile with:   pdflatex paper.tex   (twice, for references)
%  Dependencies:   amsmath, amssymb, amsthm, geometry, color, graphicx  (all standard)
%                  hyperref and T1 fontenc are used only if installed.
% ============================================================================
\documentclass[11pt,a4paper]{article}

\usepackage[utf8]{inputenc}
\usepackage{amsmath,amssymb,amsthm}
\usepackage[margin=3.2cm]{geometry}
% hyperref needs the oberdiek bundle (texlive-latexrecommended); optional.
\IfFileExists{infwarerr.sty}{\usepackage[hidelinks]{hyperref}}{}
\providecommand{\texorpdfstring}[2]{#1}   % no-op when hyperref is absent

% T1 encoding needs the ec fonts (texlive-fontsrecommended); use it only if present.
\IfFileExists{ecrm1000.tfm}{\usepackage[T1]{fontenc}}{}

% Both figures are drawn in the base LaTeX picture environment, so nothing beyond
% the standard graphics bundle is needed -- no TikZ/PGF, no external image files.
\usepackage{color,graphicx}
\definecolor{azurite} {rgb}{0.102,0.396,0.596}  % 26,101,152
\definecolor{oxide}   {rgb}{0.635,0.302,0.169}  % 162,77,43
\definecolor{cellpale}{rgb}{0.804,0.867,0.910}  % azurite, 22% on white
\definecolor{gridgray}{rgb}{0.702,0.702,0.702}
\definecolor{rulegray}{rgb}{0.502,0.502,0.502}

\emergencystretch=1.2em

\theoremstyle{plain}
\newtheorem{theorem}{Theorem}[section]
\newtheorem{proposition}[theorem]{Proposition}
\newtheorem{lemma}[theorem]{Lemma}
\newtheorem{corollary}[theorem]{Corollary}
\newtheorem{conjecture}[theorem]{Conjecture}
\theoremstyle{definition}
\newtheorem{definition}[theorem]{Definition}
\newtheorem{hypothesis}[theorem]{Hypothesis}
\newtheorem{example}[theorem]{Example}
\theoremstyle{remark}
\newtheorem{remark}[theorem]{Remark}

\newcommand{\Z}{\mathbb{Z}}
\newcommand{\Q}{\mathbb{Q}}
\newcommand{\R}{\mathbb{R}}
\newcommand{\Zn}{\Z/n\Z}
\newcommand{\units}[1]{(\Z/#1\Z)^{\times}}
\newcommand{\qint}[1]{[#1]_q}
\newcommand{\qbin}[2]{\genfrac[]{0pt}{}{#1}{#2}_{q}}
\newcommand{\qbinb}[3]{\genfrac[]{0pt}{}{#1}{#2}_{#3}}
\newcommand{\Fix}{\operatorname{Fix}}
\newcommand{\rad}{\operatorname{rad}}
\newcommand{\ord}{\operatorname{ord}}
\newcommand{\CSP}{\textup{CSP}}

\title{A crystallographic restriction for multiplier\\ cyclic sieving on necklaces}

\author{Robin Rovenszky\thanks{Discovered, verified and drafted in collaboration with
Claude (Anthropic). All computations are exact and reproducible from the accompanying
source code; see Section~\ref{sec:verification}.}}

\date{August 2026}

\begin{document}
\maketitle

\begin{abstract}
Let $Y_{n,k}(q)=\qbin{n}{k}\big/\qint{n}$ be the natural $q$-analogue of the number
$\binom{n}{k}/n$ of binary necklaces of length $n$ with $k$ black beads, defined whenever
$\gcd(n,k)=1$. Rotation acts trivially on necklaces, but the \emph{multipliers}
$x\mapsto vx$ for $v\in\units{n}$ do not; group-theoretically these are the symmetries
lying in the normaliser of the rotation group inside $S_n$. We ask when
$Y_{n,k}$ evaluated at a root of unity counts the necklaces a multiplier fixes.

Writing $f_e=\gcd(v^e-1,n)$ for the number of points of $\Zn$ held fixed by $v^e$, the
criterion reads: the sieving identity holds for every $k$ coprime to $n$ precisely when, for
every proper divisor $e$ of the order $m$ of $v$, one has $\varphi(f_e)\le 2$ --- that is
$f_e\in\{1,2,3,4,6\}$ --- together with $\rad(f_e)\mid m$. We prove that these conditions
are \emph{sufficient}, for every multiplier and with no uniformity hypothesis
(Theorem~\ref{thm:sufficiency}).

Three ingredients make this possible. First, an unconditional collapse: every $v$-fixed
necklace carries exactly $f_1$ representatives fixed by $v$ as words, so that counting fixed
necklaces reduces to the cycle index of a single permutation (Theorem~\ref{thm:collapse});
this holds for every colour content, not only for two colours. Second, a structure theorem:
under the criterion the points of $\Zn$ whose $v$-orbit is shorter than $m$ form a cyclic
subgroup $C_L$ with $L\in\{1,2,3,4,6\}$ on which $v$ acts as $\pm1$
(Theorem~\ref{thm:structure}). Third, the observation that the orbit polynomial
$\prod_{|O|<m}(1+z^{|O|})$ is palindromic of degree $L$ with linear coefficient $f_1$, while
the residues $k\bmod m$ attainable subject to $\gcd(n,k)=1$ are exactly $\{1,L-1\}$ --- a
statement equivalent to $\varphi(L)\le2$. The crystallographic restriction is thus not an
input to the argument but a reformulation of its hypothesis.

We also correct the record. The irrationality hypothesis on which the converse direction was
previously made to rest is false: the smallest counterexample is $m=9$, $r=8$, $s=4$, where
the relevant ratio equals $-1$, and there are $67$ with $m\le60$, all at $r=m-1$ and all of
modulus~$1$. We prove the congruence $m\mid r-1-s(r-s)$ that rationality forces, evaluate
the ratio in closed form at $r=m-1$, and repair the necessity argument accordingly. The
first multiplier at which the gap is felt is $n=99$, $v=19$ --- just beyond the range
$n\le96$ over which the criterion had been checked --- where $Y_{99,13}(\zeta_{10})=-9$
while the true count is $84$. All computations are exact. The cases $f\in\{3,4,6\}$ lie in a
regime that Stucky \cite{stucky} explicitly leaves open.
\end{abstract}

\section{Introduction}

Looping a string makes a cyclic group act on it, and the characters of a cyclic group are
roots of unity. The sharpest modern expression of that principle is the \emph{cyclic
sieving phenomenon} of Reiner, Stanton and White \cite{rsw}; see \cite{sagan} for a survey.

\begin{definition}[\cite{rsw}]\label{def:csp}
Let $X$ be a finite set carrying an action of a cyclic group $C=\langle c\rangle$ of order
$N$, and let $X(q)\in\Z[q]$ satisfy $X(1)=|X|$. The triple $(X,X(q),C)$ exhibits the
\emph{cyclic sieving phenomenon} if for every $d\in\Z$,
\[
  \bigl|\{x\in X: c^{\,d}x=x\}\bigr| \;=\; X(q)\big|_{q=\zeta_N^{\,d}},
\]
where $\zeta_N=e^{2\pi i/N}$.
\end{definition}

Fix $n\ge 3$ and colour the positions $\Zn$ with two colours, $k$ of them black. Modulo the
rotation action of $\Zn$ these are the \emph{necklaces}; when $\gcd(n,k)=1$ the rotation
action is free and there are exactly $\binom{n}{k}/n$ of them. The natural $q$-analogue of
that count is
\begin{equation}\label{eq:Y}
  Y_{n,k}(q)\;=\;\frac{\qbin{n}{k}}{\qint{n}},
  \qquad \qint{n}=1+q+\dots+q^{\,n-1},
\end{equation}
which is a polynomial with non-negative integer coefficients precisely when $\gcd(n,k)=1$
(Proposition~\ref{prop:factor}).

Rotation itself acts trivially on necklaces, so it carries no information. What survives is
the action of the \emph{multiplier group} $\units{n}$: for a unit $v$, relabel positions by
$x\mapsto vx$. Inside $S_n$ these are exactly the permutations normalising the cyclic
rotation group, so the full symmetry available is the affine group
$\Zn\rtimes\units{n}$. Rotation spins a necklace; a multiplier rewires it.

\paragraph{Known results.}
For $n$ prime the sieving always holds; this is Corollary~3.3 of Stucky \cite{stucky},
whose general theorem covers multipliers acting on $\Zn$ with exactly one or two fixed
points, subject to conditions on the associated Young subgroup. Beyond that the paper
stops, and says so:
\begin{quote}
\emph{``\dots when $n\not\equiv 1,2 \bmod m$ the situation appears much more delicate, and we
do not have a general conjecture.''} \cite[\S3]{stucky}
\end{quote}
That regime --- composite $n$, multipliers with three or more fixed points --- is the
subject of this note.

\paragraph{Results.}
Section~\ref{sec:count} proves the collapse theorem: the number of multiplier-fixed
necklaces is the number of multiplier-fixed \emph{words} divided by $f_1$
(Theorem~\ref{thm:collapse}). It holds for every colour content and under no hypotheses,
and it is what makes the question decidable in $O(n^2)$ arithmetic operations rather than by
exponential enumeration. Section~\ref{sec:eval} records the cyclotomic factorisation of
$Y_{n,k}$ and its evaluation at roots of unity, in both the case $m\nmid n$ and the
degenerate case $m\mid n$. Section~\ref{sec:rational} settles which of these evaluations are
rational; the hypothesis previously assumed here is disproved.
Section~\ref{sec:criterion} states the criterion, Section~\ref{sec:structure} reformulates
it as a statement about a subgroup, and Section~\ref{sec:sufficiency} proves sufficiency.
Section~\ref{sec:crystal} treats the converse: the fixed-point count is forced into the
crystallographic set $\{1,2,3,4,6\}$. Section~\ref{sec:verification} reports the
computational evidence and Section~\ref{sec:open} the boundaries. We are explicit
throughout about which statements are proved and which are only verified.

\section{Notation}\label{sec:prelim}

Throughout $n\ge 3$, and $\alpha=(\alpha_1,\dots,\alpha_\ell)$ is a \emph{content}: a
composition of $n$ recording how many positions receive each of $\ell$ colours. Words of
content $\alpha$ are maps $w:\Zn\to\{1,\dots,\ell\}$ with $|w^{-1}(i)|=\alpha_i$. The rotation
group $\Zn$ acts on them, freely if and only if $\gcd(\alpha)=1$; the orbits are
\emph{$\alpha$-necklaces}. The binary case is $\alpha=(k,n-k)$ with $\gcd(n,k)=1$.

For $v\in\units{n}$ we write $m=\ord_n(v)$ and let $v$ act on necklaces by
$[w]\mapsto[w\circ v]$. We put
\[
  f_e=\gcd(v^e-1,\;n)=\bigl|\{x\in\Zn: v^e x=x\}\bigr|,\qquad f=f_1 .
\]
$\Phi_d$ denotes the $d$-th cyclotomic polynomial, $\zeta_m$ a primitive $m$-th root of
unity, $\varphi$ Euler's totient, and $\rad(x)=\prod_{p\mid x}p$ the radical. We call $v$
\emph{uniform} if every $\langle v\rangle$-orbit on $\Zn$ has size $1$ or $m$.

\section{Counting multiplier-fixed necklaces}\label{sec:count}

For $c\in\Zn$ let $A_c:\Zn\to\Zn$ be the affine map $A_c(y)=vy+c$; these are exactly the
affine permutations with linear part $v$. Set
\[
  S_j \;=\; 1+v+v^2+\dots+v^{\,j-1}\ \in\ \Zn ,
\]
let $N_\alpha(A_c)$ denote the number of words of content $\alpha$ constant on the cycles of
$A_c$, and abbreviate $N_\alpha(v)=N_\alpha(A_0)$: the number of words of content $\alpha$
constant on the $\langle v\rangle$-orbits of $\Zn$, equivalently the number of words $w$
with $w\circ v=w$.

\begin{lemma}\label{lem:disjoint}
Let $\gcd(\alpha)=1$. A necklace $[w]$ of content $\alpha$ is fixed by the multiplier $v$ if
and only if $w$ is constant on the cycles of $A_c$ for some $c\in\Zn$, and that $c$ is then
unique.
\end{lemma}

\begin{proof}
\emph{Existence.} $[w]$ is fixed by $v$ iff $w\circ v=w\circ\rho_t$ for some rotation
$\rho_t(x)=x+t$. Substituting $y=x+t$ turns this into $w\circ A=w$ for the affine map
$A(y)=v(y-t)$, whose linear part is $v$. Conversely any $w$ with $w\circ A_c=w$ has
$v$-fixed necklace class.

\emph{Uniqueness.} If $w$ is constant on the cycles of both $A_c$ and $A_{c'}$ with $c\ne
c'$, then $w$ is invariant under $A_cA_{c'}^{-1}$, which is the rotation by $c-c'\ne0$.
Since $\gcd(\alpha)=1$ the rotation action on words of content $\alpha$ is free, a
contradiction.
\end{proof}

The next theorem is the engine of the paper. It says that only the indices $c\in(v-1)\Zn$
occur in Lemma~\ref{lem:disjoint}, and that the resulting correspondence between fixed words
and fixed necklaces is uniformly $f_1$-to-one. Counting fixed necklaces therefore reduces to
the cycle index of the single permutation $y\mapsto vy$. No hypothesis on $v$, on $n$ or on
the number of colours is required.

\begin{theorem}[collapse]\label{thm:collapse}
Let $\alpha$ be any content with $\gcd(\alpha)=1$ and let $v\in\units{n}$. Then:
\begin{enumerate}
\item[\textup{(i)}] $N_\alpha(A_c)=0$ unless $f_1\mid c$, that is, unless $c\in(v-1)\Zn$;
\item[\textup{(ii)}] the map $w\mapsto[w]$ from $\{w:w\circ v=w\}$ to the set of $v$-fixed
      $\alpha$-necklaces is surjective, with every fibre of size exactly $f_1$;
\item[\textup{(iii)}] $\displaystyle\#\Fix(v)\;=\;\frac{N_\alpha(v)}{f_1}$.
\end{enumerate}
\end{theorem}

\begin{proof}
(i) Let $\ell$ be a cycle length of $A_c$, say through $y$. Then
$A_c^{\,\ell}(y)=v^{\ell}y+cS_\ell=y$, so the congruence $(v^{\ell}-1)y\equiv-cS_\ell$ is
solvable and hence $f_\ell=\gcd(v^{\ell}-1,n)$ divides $cS_\ell$. Since $v-1$ divides
$v^{\ell}-1$ we have $f_1\mid f_\ell$, so $f_1\mid cS_\ell$; and $v\equiv1\pmod{f_1}$ gives
$S_\ell\equiv\ell\pmod{f_1}$. Therefore
\begin{equation}\label{eq:divide}
  f_1\;\bigm|\;c\,\ell\qquad\text{for every cycle length }\ell\text{ of }A_c,
\end{equation}
and consequently $e_c:=f_1/\gcd(f_1,c)$ divides \emph{every} cycle length of $A_c$. Suppose
$f_1\nmid c$, so that $e_c>1$, and choose a prime $p\mid e_c$. Then $p\mid f_1\mid n$, and
$p$ divides every cycle length of $A_c$. A word constant on the cycles of $A_c$ has each
colour class a union of cycles, so $p\mid\alpha_i$ for every $i$ and hence
$p\mid\gcd(\alpha)=1$, which is absurd. Thus $N_\alpha(A_c)=0$. Finally
$\{c:f_1\mid c\}=f_1\Zn$ is the image of multiplication by $v-1$, since
$\gcd(v-1,n)=f_1$.

(ii) \emph{Surjectivity.} Let $[w]$ be $v$-fixed. By Lemma~\ref{lem:disjoint} there is a $c$
with $w\circ A_c=w$, and by (i) we may write $c=(v-1)d$. Put $w'=w\circ\rho_{-d}$. Then
$w'(vx)=w(vx-d)$ and $w'(x)=w(x-d)$, and substituting $y=x-d$ shows that $w'\circ v=w'$ is
equivalent to $w(vy+(v-1)d)=w(y)$, i.e.\ to $w\circ A_c=w$. So $w'\in[w]$ is a $v$-fixed
word.

\emph{Fibres.} Let $w\circ v=w$ and $t\in\Zn$. The same substitution gives
$(w\circ\rho_t)\circ v=w\circ\rho_t$ if and only if $w\circ A_{t(1-v)}=w$, which by the
uniqueness clause of Lemma~\ref{lem:disjoint} forces $t(1-v)=0$, i.e.\ $t\in\Fix(v)$, a set
of size $f_1$. As the rotation action is free, these $f_1$ rotates are distinct words, and
they are precisely the $v$-fixed words in $[w]$.

(iii) Immediate from (ii), since $|\{w:w\circ v=w\}|=N_\alpha(v)$.
\end{proof}

\begin{corollary}\label{cor:twocolour}
Let $\gcd(n,k)=1$ and let $n_d$ denote the number of $\langle v\rangle$-orbits on $\Zn$ of
size exactly $d$. Then
\begin{equation}\label{eq:fixgf}
  \#\Fix(v)\;=\;\frac1{f_1}\,[z^k]\prod_{d\mid m}\bigl(1+z^{d}\bigr)^{n_d}.
\end{equation}
\end{corollary}

\begin{proof}
A word of content $(k,n-k)$ constant on the $\langle v\rangle$-orbits is exactly a choice of
a sub-multiset of orbits of total size $k$, and the generating function for the orbit sizes
is $\prod_{d\mid m}(1+z^d)^{n_d}$. Apply Theorem~\ref{thm:collapse}(iii).
\end{proof}

\begin{remark}
Theorem~\ref{thm:collapse} is strictly stronger than the observation that only $c$ with
$cS_m=0$ can contribute: $(v-1)\Zn$ is contained in $\{c:cS_m=0\}$, usually properly, and
the theorem also identifies the common value of the surviving terms. Computationally
\eqref{eq:fixgf} replaces a sum of cycle-index computations by a single $O(n)$ orbit
decomposition followed by one coefficient extraction.
\end{remark}

\section{The \texorpdfstring{$q$}{q}-analogue at a root of unity}\label{sec:eval}

\begin{proposition}\label{prop:factor}
Let $\gcd(n,k)=1$. Then $Y_{n,k}(q)$ of \eqref{eq:Y} is a polynomial, and
\[
  Y_{n,k}(q)\;=\;\prod_{d\in E(n,k)}\Phi_d(q),
  \quad
  E(n,k)=\bigl\{\,d\ge 2\ :\ d\nmid n \ \text{ and }\ (k\bmod d)>(n\bmod d)\,\bigr\}.
\]
\end{proposition}

\begin{proof}
From $\qint{j}=\prod_{d\mid j,\;d\ge2}\Phi_d(q)$ we get
$\qint{n}!=\prod_{d\ge2}\Phi_d^{\lfloor n/d\rfloor}$, hence
$\qbin{n}{k}=\prod_{d\ge2}\Phi_d^{\varepsilon_d}$ with
$\varepsilon_d=\lfloor n/d\rfloor-\lfloor k/d\rfloor-\lfloor (n-k)/d\rfloor$.
Writing $a=k\bmod d$, $b=(n-k)\bmod d$, $c=n\bmod d$, we have $a+b\in\{c,c+d\}$ and
$\varepsilon_d=(a+b-c)/d\in\{0,1\}$, with $\varepsilon_d=1$ iff $b=c-a+d$, i.e.\ iff $a>c$.
Also $\qint{n}=\prod_{d\mid n,\,d\ge2}\Phi_d$. For $d\mid n$ with $d\ge2$ we have
$c=0$ and, since $\gcd(n,k)=1$ forces $d\nmid k$, also $a>0=c$; so $\varepsilon_d=1$ and
every factor of $\qint{n}$ is cancelled exactly. The quotient is the stated product.
\end{proof}

\begin{proposition}[$q$-Lucas]\label{prop:lucas}
Let $\zeta$ be a primitive $m$-th root of unity, and write $n=am+r$, $k=bm+s$ with
$0\le r,s<m$. Then
\[
  \qbin{n}{k}\Big|_{q=\zeta}=\binom{a}{b}\,\qbin{r}{s}\Big|_{q=\zeta},
  \qquad
  \qint{n}\Big|_{q=\zeta}=[r]_{\zeta}.
\]
Consequently, if $m\nmid n$,
\begin{equation}\label{eq:eval}
  Y_{n,k}(\zeta)\;=\;\binom{a}{b}\cdot\frac{\qbin{r}{s}\big|_{q=\zeta}}{[r]_\zeta}.
\end{equation}
\end{proposition}

\begin{proof}
The first identity is the $q$-Lucas theorem. The second holds because
$\sum_{i=0}^{n-1}\zeta^{\,i}$ splits into $a$ complete cycles, each summing to $0$ when
$m>1$, leaving $\sum_{i=0}^{r-1}\zeta^{\,i}$.
\end{proof}

\begin{corollary}\label{cor:rational}
With $r\ge1$ and $0\le s$, the ratio in \eqref{eq:eval} satisfies
\[
  \frac{\qbin{r}{s}\big|_{q=\zeta}}{[r]_\zeta}
  =\begin{cases}
    1, & s\in\{1,\;r-1\},\\[2pt]
    0, & s>r,\\[2pt]
    1/[r]_\zeta, & s\in\{0,\,r\}.
  \end{cases}
\]
\end{corollary}

\begin{proof}
$\qbin{r}{1}=\qbin{r}{r-1}=\qint{r}$; $\qbin{r}{0}=\qbin{r}{r}=1$; and $\qbin{r}{s}=0$ for
$s>r$.
\end{proof}

Thus \eqref{eq:eval} is manifestly a rational integer for $s\in\{1,r-1\}$ and vanishes for
$s>r$. We call the remaining residues
\[
  B(r)\;=\;\{0,2,3,\dots,r-2,r\}
\]
the \emph{bad} residues. For these the ratio is an algebraic number in $\Q(\zeta)$ whose
rationality is a delicate matter, taken up in Section~\ref{sec:rational}.

Proposition~\ref{prop:lucas} degenerates when $m\mid n$, for then $r=0$ and $[n]_\zeta=0$.
That case is governed by the following variant, which will be needed in
Section~\ref{sec:sufficiency}.

\begin{proposition}\label{prop:mdivn}
Suppose $m\mid n$ and $\gcd(n,k)=1$, and set $a=n/m$, $b=\lfloor k/m\rfloor$,
$s=k\bmod m$. Then $s\ge1$ and
\begin{equation}\label{eq:evalm}
  Y_{n,k}(\zeta_m)\;=\;\binom{a-1}{b}\,\sigma_s,
  \qquad
  \sigma_s\;=\;\frac{(-1)^{s-1}\,\zeta^{-s(s-1)/2}}{[s]_\zeta}.
\end{equation}
\end{proposition}

\begin{proof}
From $\qbin{n}{k}=(\qint{n}/\qint{k})\,\qbin{n-1}{k-1}$ we obtain the polynomial identity
$Y_{n,k}(q)=\qbin{n-1}{k-1}/\qint{k}$, which avoids the vanishing of $\qint{n}$ at $\zeta$.
Since $m\mid n$ and $\gcd(n,k)=1$ we have $\gcd(k,m)=1$, so $s\ge1$ and
$[k]_\zeta=(1-\zeta^{k})/(1-\zeta)=(1-\zeta^{s})/(1-\zeta)=[s]_\zeta\ne0$. Writing
$n-1=m(a-1)+(m-1)$ and $k-1=mb+(s-1)$, Proposition~\ref{prop:lucas} gives
$\qbin{n-1}{k-1}\big|_\zeta=\binom{a-1}{b}\,\qbin{m-1}{s-1}\big|_\zeta$. Finally, for
$0\le j\le m-1$,
\begin{equation}\label{eq:qbinlast}
  \qbin{m-1}{j}\Big|_{q=\zeta}
  =\prod_{i=1}^{j}\frac{1-\zeta^{\,m-1-j+i}}{1-\zeta^{\,i}}
  =\prod_{u=1}^{j}\frac{1-\zeta^{-u}}{1-\zeta^{\,u}}
  =(-1)^{j}\,\zeta^{-j(j+1)/2},
\end{equation}
the middle equality by the substitution $u=j+1-i$ (legitimate since
$m-1-j+i\equiv-(j+1-i)\bmod m$) and the last by
$1-\zeta^{-u}=-\zeta^{-u}(1-\zeta^{u})$. Take $j=s-1$.
\end{proof}

\section{Rationality of the ratio}\label{sec:rational}

Fix a primitive $m$-th root of unity $\zeta$ and integers $0\le s\le r<m$, and put
\[
  \rho_{r,s}\;=\;\frac{\qbin{r}{s}\big|_{q=\zeta}}{[r]_\zeta}\;\in\;\Q(\zeta).
\]
This is defined and non-zero: $[r]_\zeta\ne0$ since $m\nmid r$ for $1\le r<m$, and
$\qbin{r}{s}\big|_\zeta\ne0$ since, by the computation in the proof of
Proposition~\ref{prop:factor}, the factor $\Phi_m$ divides $\qbin{r}{s}$ only when
$s\bmod m>r\bmod m$, that is only when $s>r$.

\begin{proposition}[conjugation constraint]\label{prop:conj}
If $\rho_{r,s}\in\Q$ then $m\mid r-1-s(r-s)$.
\end{proposition}

\begin{proof}
Complex conjugation on $\Q(\zeta)$ sends $\zeta\mapsto\zeta^{-1}$. The Gaussian binomial is
palindromic of degree $s(r-s)$, so $\qbinb{r}{s}{q^{-1}}=q^{-s(r-s)}\qbinb{r}{s}{q}$, while
$[r]_{q^{-1}}=(1-q^{-r})/(1-q^{-1})=q^{\,1-r}[r]_{q}$. Hence
\[
  \overline{\rho_{r,s}}\;=\;\zeta^{\,r-1-s(r-s)}\,\rho_{r,s}.
\]
A rational number equals its conjugate, and $\rho_{r,s}\ne0$ may be cancelled, leaving
$\zeta^{\,r-1-s(r-s)}=1$.
\end{proof}

\begin{corollary}\label{cor:qintirr}
For $2\le r<m$ the $q$-integer $[r]_\zeta$ is irrational. Hence
$\rho_{r,0}=\rho_{r,r}=1/[r]_\zeta$ is irrational for $2\le r<m$.
\end{corollary}

\begin{proof}
Apply Proposition~\ref{prop:conj} with $s=0$: rationality of $1/[r]_\zeta$, equivalently of
$[r]_\zeta$, would force $m\mid r-1$, impossible for $2\le r<m$.
\end{proof}

The case $r=m-1$ is completely explicit, and it is exactly where rationality does occur.

\begin{proposition}\label{prop:rlast}
For $0\le s\le m-1$,
\[
  \rho_{m-1,s}\;=\;(-1)^{s+1}\,\zeta^{\,1-s(s+1)/2},
\]
an algebraic number of modulus $1$. It is rational if and only if
$\zeta^{\,1-s(s+1)/2}=\pm1$, and its value is then $\pm1$.
\end{proposition}

\begin{proof}
By \eqref{eq:qbinlast}, $\qbin{m-1}{s}\big|_\zeta=(-1)^{s}\zeta^{-s(s+1)/2}$, while
$[m-1]_\zeta=(1-\zeta^{-1})/(1-\zeta)=-\zeta^{-1}$. Divide.
\end{proof}

\begin{remark}[the irrationality hypothesis is false]\label{rem:hypfalse}
An earlier version of this work assumed that $\rho_{r,s}$ is irrational for every $2\le r<m$
and every bad $s\in B(r)$. Proposition~\ref{prop:rlast} refutes this. The smallest
counterexample is $m=9$, $r=8$, $s=4$, where $1-s(s+1)/2=-9\equiv0\pmod9$ and hence
$\rho_{8,4}=(-1)^{5}=-1$. An exhaustive search --- made feasible by
Proposition~\ref{prop:conj}, which confines the candidates to a congruence class --- finds
exactly $67$ pairs with $m\le60$ and $s\in B(r)$ for which $\rho_{r,s}$ is rational. Every
one has $r=m-1$, and every value is $\pm1$; the first few are listed in
Table~\ref{tab:counterex}. The consequences for the converse direction are dealt with in
Section~\ref{sec:crystal}.
\end{remark}

\begin{table}[h]
\centering
\small
\begin{tabular}{rrlr}
$m$ & $r$ & bad $s$ with $\rho_{r,s}\in\Q$ & $\rho_{r,s}$\\\hline
$9$  & $8$  & $4$      & $-1$\\
$10$ & $9$  & $3,\ 6$  & $-1$\\
$14$ & $13$ & $5,\ 8$  & $+1$\\
$18$ & $17$ & $4,\ 13$ & $+1$\\
$18$ & $17$ & $7,\ 10$ & $-1$\\
$20$ & $19$ & $6,\ 13$ & $-1$\\
\end{tabular}
\caption{The first counterexamples to the irrationality hypothesis: bad residues at which
$\rho_{r,s}$ is nevertheless rational. All of them have $r=m-1$.}
\label{tab:counterex}
\end{table}

What survives is the following restricted statement, which is what
Section~\ref{sec:crystal} actually requires. It has been verified for all $m\le60$ and all
admissible $r,s$, using Proposition~\ref{prop:conj} to prune the search.

\begin{conjecture}\label{conj:irr}
Let $\zeta$ be a primitive $m$-th root of unity and $2\le r\le m-2$. Then $\rho_{r,s}$ is
irrational for every bad residue $s\in B(r)$.
\end{conjecture}

\begin{remark}\label{rem:modulusone}
A weaker statement suffices for Theorem~\ref{thm:necessity}, namely: \emph{if $\rho_{r,s}$
is rational then $|\rho_{r,s}|=1$}. Proposition~\ref{prop:rlast} proves this for $r=m-1$,
and it holds in every case computed.
\end{remark}

\section{The criterion}\label{sec:criterion}

\begin{definition}\label{def:criterion}
Let $n\ge3$ and $v\in\units{n}$ of order $m>1$, and put $f_e=\gcd(v^e-1,n)$. We say that $v$
satisfies the \emph{criterion} if for every proper divisor $e$ of $m$,
\begin{enumerate}
  \item[\textup{(i)}] $\varphi(f_e)\le 2$, equivalently $f_e\in\{1,2,3,4,6\}$; and
  \item[\textup{(ii)}] $\rad(f_e)\mid m$.
\end{enumerate}
We say that $v$ \emph{sieves} if
\[
  Y_{n,k}(\zeta_m)\;=\;\#\{\,k\text{-subset necklaces of }\Zn\text{ fixed by }v\,\}
  \qquad\text{for every }k\text{ coprime to }n .
\]
\end{definition}

Since $f_e\mid f_{e'}$ whenever $e\mid e'$, the criterion need only be tested at the maximal
proper divisors of $m$.

\begin{lemma}\label{lem:maximal}
$v$ satisfies the criterion if and only if $g_p:=f_{m/p}\in\{1,2,3,4,6\}$ and
$\rad(g_p)\mid m$ for every prime $p\mid m$.
\end{lemma}

\begin{proof}
Necessity is immediate. Conversely, every proper divisor $e$ of $m$ divides $m/p$ for some
prime $p\mid m$, so $f_e\mid g_p$; every divisor of an element of $\{1,2,3,4,6\}$ again lies
in $\{1,2,3,4,6\}$, and $\rad(f_e)\mid\rad(g_p)\mid m$.
\end{proof}

The two implications are in very different states, and we separate them.

\begin{theorem}[sufficiency]\label{thm:sufficiency}
If $v$ satisfies the criterion then $v$ sieves.
\end{theorem}

Theorem~\ref{thm:sufficiency} is proved in Section~\ref{sec:sufficiency}. It applies to
every multiplier, with no uniformity hypothesis. The converse is established in
Section~\ref{sec:crystal} for uniform multipliers with $f<m$, modulo
Conjecture~\ref{conj:irr}; in general it is verified but unproved.

\begin{conjecture}[converse]\label{conj:converse}
If $v$ sieves then $v$ satisfies the criterion.
\end{conjecture}

Since the cyclic sieving phenomenon of Definition~\ref{def:csp} for the group
$\langle u\rangle$ requires the identity at every power $u^{d}$, and $u^{d}$ is itself a
multiplier, Theorem~\ref{thm:sufficiency} determines the full \CSP{} by conjunction over the
divisors of $\ord_n(u)$.

\begin{corollary}[uniform case]\label{cor:uniform}
If $v$ is uniform then $f_e=f$ for every proper divisor $e$ of $m$, and the criterion reads:
$f\in\{1,2,3,4,6\}$ and $\rad(f)\mid m$. In particular the verdict depends only on the pair
$(m,f)$ --- not on $n$, and not on the number of orbits.
\end{corollary}

\begin{proof}
If $e\mid m$ is proper and $v^{e}x=x$, the $\langle v\rangle$-orbit of $x$ has size dividing
$e<m$, so by uniformity it has size $1$. Hence $\Fix(v^{e})=\Fix(v)$ and $f_e=f$.
\end{proof}

The independence from $n$ asserted in Corollary~\ref{cor:uniform} is a striking feature of
the data and was the first indication that a clean statement existed; see
Figure~\ref{fig:grid}.

\begin{figure}[t]
\centering
\setlength{\unitlength}{0.66cm}%
\begin{picture}(10.1,13.3)(-2.05,-0.1)
  % --- shaded cells: sieving holds
  \put(0,10){\textcolor{cellpale}{\rule{1\unitlength}{1\unitlength}}}
  \put(1,10){\textcolor{cellpale}{\rule{1\unitlength}{1\unitlength}}}
  \put(3,10){\textcolor{cellpale}{\rule{1\unitlength}{1\unitlength}}}
  \put(0,9){\textcolor{cellpale}{\rule{1\unitlength}{1\unitlength}}}
  \put(2,9){\textcolor{cellpale}{\rule{1\unitlength}{1\unitlength}}}
  \put(0,8){\textcolor{cellpale}{\rule{1\unitlength}{1\unitlength}}}
  \put(1,8){\textcolor{cellpale}{\rule{1\unitlength}{1\unitlength}}}
  \put(3,8){\textcolor{cellpale}{\rule{1\unitlength}{1\unitlength}}}
  \put(0,7){\textcolor{cellpale}{\rule{1\unitlength}{1\unitlength}}}
  \put(0,6){\textcolor{cellpale}{\rule{1\unitlength}{1\unitlength}}}
  \put(1,6){\textcolor{cellpale}{\rule{1\unitlength}{1\unitlength}}}
  \put(2,6){\textcolor{cellpale}{\rule{1\unitlength}{1\unitlength}}}
  \put(3,6){\textcolor{cellpale}{\rule{1\unitlength}{1\unitlength}}}
  \put(5,6){\textcolor{cellpale}{\rule{1\unitlength}{1\unitlength}}}
  \put(0,5){\textcolor{cellpale}{\rule{1\unitlength}{1\unitlength}}}
  \put(0,4){\textcolor{cellpale}{\rule{1\unitlength}{1\unitlength}}}
  \put(1,4){\textcolor{cellpale}{\rule{1\unitlength}{1\unitlength}}}
  \put(3,4){\textcolor{cellpale}{\rule{1\unitlength}{1\unitlength}}}
  \put(0,3){\textcolor{cellpale}{\rule{1\unitlength}{1\unitlength}}}
  \put(2,3){\textcolor{cellpale}{\rule{1\unitlength}{1\unitlength}}}
  \put(0,2){\textcolor{cellpale}{\rule{1\unitlength}{1\unitlength}}}
  \put(1,2){\textcolor{cellpale}{\rule{1\unitlength}{1\unitlength}}}
  \put(3,2){\textcolor{cellpale}{\rule{1\unitlength}{1\unitlength}}}
  \put(0,1){\textcolor{cellpale}{\rule{1\unitlength}{1\unitlength}}}
  \put(0,0){\textcolor{cellpale}{\rule{1\unitlength}{1\unitlength}}}
  \put(1,0){\textcolor{cellpale}{\rule{1\unitlength}{1\unitlength}}}
  \put(2,0){\textcolor{cellpale}{\rule{1\unitlength}{1\unitlength}}}
  \put(3,0){\textcolor{cellpale}{\rule{1\unitlength}{1\unitlength}}}
  \put(5,0){\textcolor{cellpale}{\rule{1\unitlength}{1\unitlength}}}
  % --- grid
  \thinlines
  \put(0,0){\textcolor{gridgray}{\line(1,0){8}}}
  \put(0,1){\textcolor{gridgray}{\line(1,0){8}}}
  \put(0,2){\textcolor{gridgray}{\line(1,0){8}}}
  \put(0,3){\textcolor{gridgray}{\line(1,0){8}}}
  \put(0,4){\textcolor{gridgray}{\line(1,0){8}}}
  \put(0,5){\textcolor{gridgray}{\line(1,0){8}}}
  \put(0,6){\textcolor{gridgray}{\line(1,0){8}}}
  \put(0,7){\textcolor{gridgray}{\line(1,0){8}}}
  \put(0,8){\textcolor{gridgray}{\line(1,0){8}}}
  \put(0,9){\textcolor{gridgray}{\line(1,0){8}}}
  \put(0,10){\textcolor{gridgray}{\line(1,0){8}}}
  \put(0,11){\textcolor{gridgray}{\line(1,0){8}}}
  \put(0,0){\textcolor{gridgray}{\line(0,1){11}}}
  \put(1,0){\textcolor{gridgray}{\line(0,1){11}}}
  \put(2,0){\textcolor{gridgray}{\line(0,1){11}}}
  \put(3,0){\textcolor{gridgray}{\line(0,1){11}}}
  \put(4,0){\textcolor{gridgray}{\line(0,1){11}}}
  \put(5,0){\textcolor{gridgray}{\line(0,1){11}}}
  \put(6,0){\textcolor{gridgray}{\line(0,1){11}}}
  \put(7,0){\textcolor{gridgray}{\line(0,1){11}}}
  \put(8,0){\textcolor{gridgray}{\line(0,1){11}}}
  % --- divider: known cases (f=1,2) | new cases
  \put(2,0){\textcolor{rulegray}{\rule{0.9pt}{11\unitlength}}}
  % --- verdict marks
  \put(0,10){\makebox(1,1){\textcolor{azurite}{$\bullet$}}}
  \put(1,10){\makebox(1,1){\textcolor{azurite}{$\bullet$}}}
  \put(3,10){\makebox(1,1){\textcolor{azurite}{$\bullet$}}}
  \put(0,9){\makebox(1,1){\textcolor{azurite}{$\bullet$}}}
  \put(2,9){\makebox(1,1){\textcolor{azurite}{$\bullet$}}}
  \put(0,8){\makebox(1,1){\textcolor{azurite}{$\bullet$}}}
  \put(1,8){\makebox(1,1){\textcolor{azurite}{$\bullet$}}}
  \put(3,8){\makebox(1,1){\textcolor{azurite}{$\bullet$}}}
  \put(0,7){\makebox(1,1){\textcolor{azurite}{$\bullet$}}}
  \put(0,6){\makebox(1,1){\textcolor{azurite}{$\bullet$}}}
  \put(1,6){\makebox(1,1){\textcolor{azurite}{$\bullet$}}}
  \put(2,6){\makebox(1,1){\textcolor{azurite}{$\bullet$}}}
  \put(3,6){\makebox(1,1){\textcolor{azurite}{$\bullet$}}}
  \put(5,6){\makebox(1,1){\textcolor{azurite}{$\bullet$}}}
  \put(0,5){\makebox(1,1){\textcolor{azurite}{$\bullet$}}}
  \put(0,4){\makebox(1,1){\textcolor{azurite}{$\bullet$}}}
  \put(1,4){\makebox(1,1){\textcolor{azurite}{$\bullet$}}}
  \put(3,4){\makebox(1,1){\textcolor{azurite}{$\bullet$}}}
  \put(0,3){\makebox(1,1){\textcolor{azurite}{$\bullet$}}}
  \put(2,3){\makebox(1,1){\textcolor{azurite}{$\bullet$}}}
  \put(0,2){\makebox(1,1){\textcolor{azurite}{$\bullet$}}}
  \put(1,2){\makebox(1,1){\textcolor{azurite}{$\bullet$}}}
  \put(3,2){\makebox(1,1){\textcolor{azurite}{$\bullet$}}}
  \put(0,1){\makebox(1,1){\textcolor{azurite}{$\bullet$}}}
  \put(0,0){\makebox(1,1){\textcolor{azurite}{$\bullet$}}}
  \put(1,0){\makebox(1,1){\textcolor{azurite}{$\bullet$}}}
  \put(2,0){\makebox(1,1){\textcolor{azurite}{$\bullet$}}}
  \put(3,0){\makebox(1,1){\textcolor{azurite}{$\bullet$}}}
  \put(5,0){\makebox(1,1){\textcolor{azurite}{$\bullet$}}}
  % --- axis ticks
  \put(0,11.13){\makebox(1,0.55)[b]{\footnotesize $1$}}
  \put(1,11.13){\makebox(1,0.55)[b]{\footnotesize $2$}}
  \put(2,11.13){\makebox(1,0.55)[b]{\footnotesize $3$}}
  \put(3,11.13){\makebox(1,0.55)[b]{\footnotesize $4$}}
  \put(4,11.13){\makebox(1,0.55)[b]{\footnotesize $5$}}
  \put(5,11.13){\makebox(1,0.55)[b]{\footnotesize $6$}}
  \put(6,11.13){\makebox(1,0.55)[b]{\footnotesize $7$}}
  \put(7,11.13){\makebox(1,0.55)[b]{\footnotesize $8$}}
  \put(-1.42,10){\makebox(1.25,1)[r]{\footnotesize $2$}}
  \put(-1.42,9){\makebox(1.25,1)[r]{\footnotesize $3$}}
  \put(-1.42,8){\makebox(1.25,1)[r]{\footnotesize $4$}}
  \put(-1.42,7){\makebox(1.25,1)[r]{\footnotesize $5$}}
  \put(-1.42,6){\makebox(1.25,1)[r]{\footnotesize $6$}}
  \put(-1.42,5){\makebox(1.25,1)[r]{\footnotesize $7$}}
  \put(-1.42,4){\makebox(1.25,1)[r]{\footnotesize $8$}}
  \put(-1.42,3){\makebox(1.25,1)[r]{\footnotesize $9$}}
  \put(-1.42,2){\makebox(1.25,1)[r]{\footnotesize $10$}}
  \put(-1.42,1){\makebox(1.25,1)[r]{\footnotesize $11$}}
  \put(-1.42,0){\makebox(1.25,1)[r]{\footnotesize $12$}}
  % --- axis titles
  \put(0,12.34){\makebox(8,0.6)[b]{\small fixed-point count $f$}}
  \put(-2.05,0){\makebox(0.6,11){\rotatebox{90}{\small order $m$}}}
\end{picture}
\caption{Uniform multipliers: sieving holds ($\bullet$, shaded) or fails (blank), by the
order $m$ of the multiplier and its fixed-point count $f$. Columns $f=5,7,8$ are empty for
every $m$; the pattern of a row depends only on which primes divide $m$. Columns $f=1,2$ are
the previously known cases, columns $f=3,4,6$ are new.}
\label{fig:grid}
\end{figure}

\begin{figure}[t]
\centering
\setlength{\unitlength}{1cm}%
\begin{picture}(13.2,2.15)(-0.66,-1.48)
  % --- f = 1 (allowed)
  \put(0,0){\textcolor{gridgray}{\circle{1.24}}}
  \put(0,0.62){\textcolor{azurite}{\circle*{0.16}}}
  \put(-0.5,-1.04){\makebox(1,0.42)[b]{\footnotesize $1$}}
  \put(-0.8,-1.44){\makebox(1.6,0.36)[b]{\textcolor{azurite}{\tiny allowed}}}
  % --- f = 2 (allowed)
  \put(1.7,0){\textcolor{gridgray}{\circle{1.24}}}
  \put(1.7,0.62){\textcolor{azurite}{\circle*{0.16}}}
  \put(1.7,-0.62){\textcolor{azurite}{\circle*{0.16}}}
  \put(1.2,-1.04){\makebox(1,0.42)[b]{\footnotesize $2$}}
  \put(0.9,-1.44){\makebox(1.6,0.36)[b]{\textcolor{azurite}{\tiny allowed}}}
  % --- f = 3 (allowed)
  \put(3.4,0){\textcolor{gridgray}{\circle{1.24}}}
  \put(3.4,0.62){\textcolor{azurite}{\circle*{0.16}}}
  \put(2.8631,-0.31){\textcolor{azurite}{\circle*{0.16}}}
  \put(3.9369,-0.31){\textcolor{azurite}{\circle*{0.16}}}
  \put(2.9,-1.04){\makebox(1,0.42)[b]{\footnotesize $3$}}
  \put(2.6,-1.44){\makebox(1.6,0.36)[b]{\textcolor{azurite}{\tiny allowed}}}
  % --- f = 4 (allowed)
  \put(5.1,0){\textcolor{gridgray}{\circle{1.24}}}
  \put(5.1,0.62){\textcolor{azurite}{\circle*{0.16}}}
  \put(4.48,0){\textcolor{azurite}{\circle*{0.16}}}
  \put(5.1,-0.62){\textcolor{azurite}{\circle*{0.16}}}
  \put(5.72,0){\textcolor{azurite}{\circle*{0.16}}}
  \put(4.6,-1.04){\makebox(1,0.42)[b]{\footnotesize $4$}}
  \put(4.3,-1.44){\makebox(1.6,0.36)[b]{\textcolor{azurite}{\tiny allowed}}}
  % --- f = 5 (forbidden)
  \put(6.8,0){\textcolor{gridgray}{\circle{1.24}}}
  \put(6.8,0.62){\textcolor{oxide}{\circle*{0.16}}}
  \put(6.2103,0.1916){\textcolor{oxide}{\circle*{0.16}}}
  \put(6.4356,-0.5016){\textcolor{oxide}{\circle*{0.16}}}
  \put(7.1644,-0.5016){\textcolor{oxide}{\circle*{0.16}}}
  \put(7.3897,0.1916){\textcolor{oxide}{\circle*{0.16}}}
  \put(6.3,-1.04){\makebox(1,0.42)[b]{\footnotesize $5$}}
  \put(6,-1.44){\makebox(1.6,0.36)[b]{\textcolor{oxide}{\tiny forbidden}}}
  % --- f = 6 (allowed)
  \put(8.5,0){\textcolor{gridgray}{\circle{1.24}}}
  \put(8.5,0.62){\textcolor{azurite}{\circle*{0.16}}}
  \put(7.9631,0.31){\textcolor{azurite}{\circle*{0.16}}}
  \put(7.9631,-0.31){\textcolor{azurite}{\circle*{0.16}}}
  \put(8.5,-0.62){\textcolor{azurite}{\circle*{0.16}}}
  \put(9.0369,-0.31){\textcolor{azurite}{\circle*{0.16}}}
  \put(9.0369,0.31){\textcolor{azurite}{\circle*{0.16}}}
  \put(8,-1.04){\makebox(1,0.42)[b]{\footnotesize $6$}}
  \put(7.7,-1.44){\makebox(1.6,0.36)[b]{\textcolor{azurite}{\tiny allowed}}}
  % --- f = 7 (forbidden)
  \put(10.2,0){\textcolor{gridgray}{\circle{1.24}}}
  \put(10.2,0.62){\textcolor{oxide}{\circle*{0.16}}}
  \put(9.7153,0.3866){\textcolor{oxide}{\circle*{0.16}}}
  \put(9.5955,-0.138){\textcolor{oxide}{\circle*{0.16}}}
  \put(9.931,-0.5586){\textcolor{oxide}{\circle*{0.16}}}
  \put(10.469,-0.5586){\textcolor{oxide}{\circle*{0.16}}}
  \put(10.8045,-0.138){\textcolor{oxide}{\circle*{0.16}}}
  \put(10.6847,0.3866){\textcolor{oxide}{\circle*{0.16}}}
  \put(9.7,-1.04){\makebox(1,0.42)[b]{\footnotesize $7$}}
  \put(9.4,-1.44){\makebox(1.6,0.36)[b]{\textcolor{oxide}{\tiny forbidden}}}
  % --- f = 8 (forbidden)
  \put(11.9,0){\textcolor{gridgray}{\circle{1.24}}}
  \put(11.9,0.62){\textcolor{oxide}{\circle*{0.16}}}
  \put(11.4616,0.4384){\textcolor{oxide}{\circle*{0.16}}}
  \put(11.28,0){\textcolor{oxide}{\circle*{0.16}}}
  \put(11.4616,-0.4384){\textcolor{oxide}{\circle*{0.16}}}
  \put(11.9,-0.62){\textcolor{oxide}{\circle*{0.16}}}
  \put(12.3384,-0.4384){\textcolor{oxide}{\circle*{0.16}}}
  \put(12.52,0){\textcolor{oxide}{\circle*{0.16}}}
  \put(12.3384,0.4384){\textcolor{oxide}{\circle*{0.16}}}
  \put(11.4,-1.04){\makebox(1,0.42)[b]{\footnotesize $8$}}
  \put(11.1,-1.44){\makebox(1.6,0.36)[b]{\textcolor{oxide}{\tiny forbidden}}}
\end{picture}
\caption{The permitted sizes of a multiplier's fixed set, $f\in\{1,2,3,4,6\}$
(Theorem~\ref{thm:necessity}) --- identical to the permitted orders of rotational symmetry
of a plane lattice, and for the same underlying reason: a quantity built from a root of
unity is required to be rational.}
\label{fig:orders}
\end{figure}

\section{What the criterion says}\label{sec:structure}

Write
\[
  P\;=\;\bigl\{\,y\in\Zn:\ \text{the }\langle v\rangle\text{-orbit of }y
        \text{ has size}<m\,\bigr\}
   \;=\;\bigcup_{p\mid m\ \mathrm{prime}}\Fix\bigl(v^{\,m/p}\bigr),
\]
the second equality because a period smaller than $m$ divides $m/p$ for some prime $p\mid m$.
Each $\Fix(v^{e})$ is the unique subgroup of the cyclic group $\Zn$ of order $f_e$, so $P$ is
a union of cyclic subgroups of orders $g_p=f_{m/p}$.

\begin{lemma}[no clash]\label{lem:noclash}
Assume $v$ satisfies the criterion. Then the integers $g_p$, for $p\mid m$ prime, are totally
ordered by divisibility.
\end{lemma}

\begin{proof}
By Lemma~\ref{lem:maximal} every $g_p$ lies in $\{1,2,3,4,6\}$, whose incomparable pairs are
$\{2,3\}$, $\{3,4\}$ and $\{4,6\}$. Two facts are used throughout: $g_p\mid n$; and if
$2\mid n$ then $v$ is odd, so $2\mid v-1\mid v^{e}-1$ and hence $2\mid f_e$ for every $e$.

\emph{$\{2,3\}$.} If $g_p=2$ then $2\mid n$, so every $f_e$ is even, contradicting the
oddness of $g_q=3$.

\emph{$\{3,4\}$.} If $g_q=4$ then $4\mid n$, so again every $f_e$ is even, contradicting
$g_p=3$.

\emph{$\{4,6\}$.} Suppose $g_p=4$ and $g_q=6$ with $p\ne q$. Then $4\mid n$ and $6\mid n$, so
$12\mid n$. Since $\gcd(v^{m/p}-1,n)=4$ and $3\mid n$, we get $3\nmid v^{m/p}-1$; hence
$\ord_3(v)\ne1$, so $\ord_3(v)=2$ and $2\nmid m/p$. Since $\gcd(v^{m/q}-1,n)=6$ and
$4\mid n$, we get $4\nmid v^{m/q}-1$; hence $\ord_4(v)=2$ and $2\nmid m/q$. But $3\mid n$
forces $\ord_3(v)\mid m$, so $2\mid m$; then $2\nmid m/p$ forces $p=2$, and likewise $q=2$,
contradicting $p\ne q$.
\end{proof}

\begin{theorem}[structure]\label{thm:structure}
Let $v\in\units{n}$ have order $m>1$ and put $L=\max_p g_p$. If $v$ satisfies the criterion
then $P$ is precisely the cyclic subgroup $C_L\le\Zn$ of order $L$, with
\[
  L\in\{1,2,3,4,6\},\qquad \rad(L)\mid m,
\]
and $v$ acts on $C_L$ as $+1$ or $-1$. Conversely, if $P$ is a subgroup of order $L$ with
$\varphi(L)\le2$ and $\rad(L)\mid m$, then $v$ satisfies the criterion.
\end{theorem}

\begin{proof}
Assume the criterion. By Lemma~\ref{lem:noclash} the subgroups $\Fix(v^{m/p})=C_{g_p}$ are
nested, so their union $P$ is the largest of them, namely $C_L$; and $L=g_{p_0}$ lies in
$\{1,2,3,4,6\}$ with $\rad(L)\mid m$ by Lemma~\ref{lem:maximal}. The subgroup $C_L$ is
$\langle v\rangle$-stable, and under the isomorphism $C_L\cong\Z/L\Z$ the multiplier acts as
multiplication by $v\bmod L$, an element of $\units{L}$. That group has order
$\varphi(L)\le2$ and therefore equals $\{\pm1\}$.

Conversely, if $P$ is a subgroup of order $L$ with $\varphi(L)\le2$ and $\rad(L)\mid m$, then
for every prime $p\mid m$ the group $\Fix(v^{m/p})$ is a subgroup of $P$, so $g_p\mid L$.
Every divisor of an element of $\{1,2,3,4,6\}$ again lies in that set, and
$\rad(g_p)\mid\rad(L)\mid m$; apply Lemma~\ref{lem:maximal}.
\end{proof}

Theorem~\ref{thm:structure} is a reformulation rather than a consequence: \emph{the criterion
says exactly that the points of period less than $m$ form a subgroup of crystallographic
order whose radical divides $m$.} Since $v$ then acts on that subgroup as $\pm1$, its orbits
have size $1$ or $2$, and only eight configurations are possible; all eight occur
(Figure~\ref{fig:shapes}).

\begin{figure}[h]
\centering
\setlength{\unitlength}{1cm}%
\begin{picture}(13.2,2.25)(-0.66,-1.58)
  % --- L=1, eps=+1
  \put(0,0){\textcolor{gridgray}{\circle{1.24}}}
  \put(0,0.62){\textcolor{azurite}{\circle*{0.16}}}
  \put(-0.6,-1.04){\makebox(1.2,0.42)[b]{\footnotesize $L=1$}}
  \put(-0.8,-1.5){\makebox(1.6,0.36)[b]{\tiny $\varepsilon=+1$}}
  % --- L=2, eps=+1
  \put(1.7,0){\textcolor{gridgray}{\circle{1.24}}}
  \put(1.7,0.62){\textcolor{azurite}{\circle*{0.16}}}
  \put(1.7,-0.62){\textcolor{azurite}{\circle*{0.16}}}
  \put(1.1,-1.04){\makebox(1.2,0.42)[b]{\footnotesize $L=2$}}
  \put(0.9,-1.5){\makebox(1.6,0.36)[b]{\tiny $\varepsilon=+1$}}
  % --- L=3, eps=+1
  \put(3.4,0){\textcolor{gridgray}{\circle{1.24}}}
  \put(3.4,0.62){\textcolor{azurite}{\circle*{0.16}}}
  \put(3.9369,-0.31){\textcolor{azurite}{\circle*{0.16}}}
  \put(2.8631,-0.31){\textcolor{azurite}{\circle*{0.16}}}
  \put(2.8,-1.04){\makebox(1.2,0.42)[b]{\footnotesize $L=3$}}
  \put(2.6,-1.5){\makebox(1.6,0.36)[b]{\tiny $\varepsilon=+1$}}
  % --- L=3, eps=-1
  \put(5.1,0){\textcolor{gridgray}{\circle{1.24}}}
  \put(4.5631,-0.31){\textcolor{oxide}{\line(1,0){1.0738}}}
  \put(5.1,0.62){\textcolor{azurite}{\circle*{0.16}}}
  \put(5.6369,-0.31){\textcolor{oxide}{\circle*{0.16}}}
  \put(4.5631,-0.31){\textcolor{oxide}{\circle*{0.16}}}
  \put(4.5,-1.04){\makebox(1.2,0.42)[b]{\footnotesize $L=3$}}
  \put(4.3,-1.5){\makebox(1.6,0.36)[b]{\tiny $\varepsilon=-1$}}
  % --- L=4, eps=+1
  \put(6.8,0){\textcolor{gridgray}{\circle{1.24}}}
  \put(6.8,0.62){\textcolor{azurite}{\circle*{0.16}}}
  \put(7.42,0){\textcolor{azurite}{\circle*{0.16}}}
  \put(6.8,-0.62){\textcolor{azurite}{\circle*{0.16}}}
  \put(6.18,0){\textcolor{azurite}{\circle*{0.16}}}
  \put(6.2,-1.04){\makebox(1.2,0.42)[b]{\footnotesize $L=4$}}
  \put(6.0,-1.5){\makebox(1.6,0.36)[b]{\tiny $\varepsilon=+1$}}
  % --- L=4, eps=-1
  \put(8.5,0){\textcolor{gridgray}{\circle{1.24}}}
  \put(7.88,0){\textcolor{oxide}{\line(1,0){1.24}}}
  \put(8.5,0.62){\textcolor{azurite}{\circle*{0.16}}}
  \put(8.5,-0.62){\textcolor{azurite}{\circle*{0.16}}}
  \put(9.12,0){\textcolor{oxide}{\circle*{0.16}}}
  \put(7.88,0){\textcolor{oxide}{\circle*{0.16}}}
  \put(7.9,-1.04){\makebox(1.2,0.42)[b]{\footnotesize $L=4$}}
  \put(7.7,-1.5){\makebox(1.6,0.36)[b]{\tiny $\varepsilon=-1$}}
  % --- L=6, eps=+1
  \put(10.2,0){\textcolor{gridgray}{\circle{1.24}}}
  \put(10.2,0.62){\textcolor{azurite}{\circle*{0.16}}}
  \put(10.7369,0.31){\textcolor{azurite}{\circle*{0.16}}}
  \put(10.7369,-0.31){\textcolor{azurite}{\circle*{0.16}}}
  \put(10.2,-0.62){\textcolor{azurite}{\circle*{0.16}}}
  \put(9.6631,-0.31){\textcolor{azurite}{\circle*{0.16}}}
  \put(9.6631,0.31){\textcolor{azurite}{\circle*{0.16}}}
  \put(9.6,-1.04){\makebox(1.2,0.42)[b]{\footnotesize $L=6$}}
  \put(9.4,-1.5){\makebox(1.6,0.36)[b]{\tiny $\varepsilon=+1$}}
  % --- L=6, eps=-1
  \put(11.9,0){\textcolor{gridgray}{\circle{1.24}}}
  \put(11.3631,0.31){\textcolor{oxide}{\line(1,0){1.0738}}}
  \put(11.3631,-0.31){\textcolor{oxide}{\line(1,0){1.0738}}}
  \put(11.9,0.62){\textcolor{azurite}{\circle*{0.16}}}
  \put(11.9,-0.62){\textcolor{azurite}{\circle*{0.16}}}
  \put(12.4369,0.31){\textcolor{oxide}{\circle*{0.16}}}
  \put(12.4369,-0.31){\textcolor{oxide}{\circle*{0.16}}}
  \put(11.3631,-0.31){\textcolor{oxide}{\circle*{0.16}}}
  \put(11.3631,0.31){\textcolor{oxide}{\circle*{0.16}}}
  \put(11.3,-1.04){\makebox(1.2,0.42)[b]{\footnotesize $L=6$}}
  \put(11.1,-1.5){\makebox(1.6,0.36)[b]{\tiny $\varepsilon=-1$}}
\end{picture}
\caption{The eight admissible shapes of the non-generic part $P=C_L$ under the criterion
(Theorem~\ref{thm:structure}). Points fixed by the multiplier are drawn in azurite; the
two-element orbits, which occur when $v$ acts as $-1$, are drawn in oxide and joined by their
chord. Every criterion-satisfying multiplier with $n\le400$ realises one of these eight, and
each of the eight is realised.}
\label{fig:shapes}
\end{figure}

\section{Proof of sufficiency}\label{sec:sufficiency}

Throughout this section $v$ satisfies the criterion, and $L$ and $C_L$ are as in
Theorem~\ref{thm:structure}. Write $f=f_1$, let $n_m$ be the number of
$\langle v\rangle$-orbits of size $m$, and set
\[
  Q(z)\;=\;\prod_{|O|<m}\bigl(1+z^{|O|}\bigr),
\]
the product taken over the orbits contained in $P$. Then $n=L+mn_m$, and
Corollary~\ref{cor:twocolour} reads
\begin{equation}\label{eq:split}
  \#\Fix(v)\;=\;\frac1{f}\,[z^k]\,N(z),
  \qquad N(z)\;=\;Q(z)\,\bigl(1+z^{m}\bigr)^{n_m}.
\end{equation}
Identity \eqref{eq:split}, and the expansion \eqref{eq:expand} below, are formal consequences
of Corollary~\ref{cor:twocolour} and the definition of $Q$; they hold for every multiplier.
Only the three lemmas that follow, and the proof of Theorem~\ref{thm:sufficiency}, invoke the
criterion.

\begin{lemma}\label{lem:Q}
$Q$ is palindromic of degree $L$, and $[z^{1}]Q=[z^{\,L-1}]Q=f$.
\end{lemma}

\begin{proof}
Each factor $1+z^{|O|}$ is palindromic, hence so is the product, and
$\deg Q=\sum_{|O|<m}|O|=|P|=L$. The coefficient of $z^{1}$ counts the orbits of size $1$,
that is $|\Fix(v)|=f$, and palindromicity gives $[z^{L-1}]Q=[z^{1}]Q$.
\end{proof}

\begin{lemma}\label{lem:residuesL}
Let $s=k\bmod m$ for some $k$ coprime to $n$. Then $\gcd(s,\rad(L))=1$. If in addition
$L\ge2$ and $0\le s\le L$, then $s\in\{1,L-1\}$.
\end{lemma}

\begin{proof}
Let $p$ be a prime dividing $\rad(L)$. Then $p\mid L\mid n$ and $p\mid\rad(L)\mid m$; if
$p\mid s$ then $p\mid k$, contradicting $\gcd(n,k)=1$. For the second claim, $s=0$ and $s=L$
are divisible by every prime factor of $L$ and so are excluded. It remains to see that every
$s$ with $2\le s\le L-2$ shares a prime factor with $L$; this range is empty for $L=2,3$,
consists of $s=2$ for $L=4$, and of $s=2,3,4$ for $L=6$, and in each case the assertion
holds. It is precisely here that $\varphi(L)\le2$ is used: for $L=5$ the residue $s=2$ would
survive, and for $L\ge7$ some prime below $L-1$ fails to divide $L$.
\end{proof}

\begin{lemma}\label{lem:dichotomy}
Put $r=n\bmod m$. Then $r=L\bmod m$, and exactly one of the following holds:
\begin{enumerate}
\item[\textup{(a)}] $m>L$; then $r=L\ge1$ and $m\nmid n$;
\item[\textup{(b)}] $m\le L$; then $(L,m)\in\{(2,2),(3,3),(4,2),(4,4),(6,6)\}$, $r=0$ and
      $m\mid n$.
\end{enumerate}
\end{lemma}

\begin{proof}
$n=L+mn_m$ gives $r=L\bmod m$. In case (a), $r=L$, and $L\ge1$ because $0\in P$. In case
(b), $L\in\{1,2,3,4,6\}$ with $\rad(L)\mid m$ and $2\le m\le L$ leaves exactly the five
listed pairs, and in each $m\mid L$, so $r=0$.
\end{proof}

\begin{proof}[Proof of Theorem~\textup{\ref{thm:sufficiency}}]
Fix $k$ coprime to $n$ and put $b=\lfloor k/m\rfloor$ and $s=k\bmod m$. By \eqref{eq:split}
it suffices to show $[z^k]N(z)=f\,Y_{n,k}(\zeta_m)$. Expanding \eqref{eq:split} and
substituting $i=b-j$ for the exponent of $(1+z^m)^{n_m}$,
\begin{equation}\label{eq:expand}
  [z^k]N\;=\;\sum_{j\ge0}\binom{n_m}{\,b-j\,}\,[z^{\,jm+s}]\,Q .
\end{equation}

\emph{Case \textup{(a)}: $m\nmid n$.} By Lemma~\ref{lem:dichotomy}, $r=L<m$, and
$n_m=\lfloor n/m\rfloor=:a$ because $n=L+mn_m$ with $L<m$. Every
term of \eqref{eq:expand} with $j\ge1$ vanishes, since then $jm+s\ge m>L=\deg Q$; hence
$[z^k]N=\binom{a}{b}\,[z^{s}]Q$. On the other side \eqref{eq:eval} gives
$Y_{n,k}(\zeta_m)=\binom{a}{b}\rho_{L,s}$, so it suffices to prove $[z^{s}]Q=f\rho_{L,s}$.
If $s>L$ then both sides vanish, because $\deg Q=L$ and $\qbin{L}{s}=0$. If $L\ge2$ and
$s\le L$ then $s\in\{1,L-1\}$ by Lemma~\ref{lem:residuesL}, whence $\rho_{L,s}=1$ by
Corollary~\ref{cor:rational}, while $[z^{s}]Q=f$ by Lemma~\ref{lem:Q}. If $L=1$ then $f=1$, $Q=1+z$, and $\rho_{1,s}=1$ for $s\in\{0,1\}$, so
both sides equal $1$.

\emph{Case \textup{(b)}: $m\mid n$.} Here $r=0$, $u:=L/m\ge1$ is an integer and
$a:=n/m=n_m+u$. By Vandermonde's identity
$\binom{a-1}{b}=\sum_{j\ge0}\binom{u-1}{j}\binom{n_m}{\,b-j\,}$, so comparing
\eqref{eq:expand} with \eqref{eq:evalm} it suffices to prove
\begin{equation}\label{eq:target}
  [z^{\,jm+s}]Q\;=\;f\,\sigma_s\binom{u-1}{j}\qquad\text{for every }j\ge0 .
\end{equation}
For each of the five pairs $\rad(L)=\rad(m)$, so Lemma~\ref{lem:residuesL} says that the
attainable residues are those $s\in\{0,\dots,m-1\}$ coprime to $m$, namely $s\in\{1,m-1\}$.
Now $\sigma_1=1$, and by \eqref{eq:qbinlast} together with $[m-1]_\zeta=-\zeta^{-1}$,
\[
  \sigma_{m-1}\;=\;(-1)^{m-1}\,\zeta^{\,1-(m-1)(m-2)/2}\;=\;1
  \qquad\text{for } m\in\{2,3,4,6\},
\]
the four values of the exponent being $1,0,-2,-9$ against $\zeta$ of order $2,3,4,6$
respectively. Hence \eqref{eq:target} reduces to $[z^{\,jm+s}]Q=f\binom{u-1}{j}$.

If $u=1$, that is $(L,m)\in\{(2,2),(3,3),(4,4),(6,6)\}$, then $\binom{0}{j}$ vanishes unless
$j=0$. For $j=0$ we need $[z^{s}]Q=f$ with $s\in\{1,m-1\}=\{1,L-1\}$, which is
Lemma~\ref{lem:Q}; for $j\ge1$ we have $jm+s>m=L=\deg Q$ and both sides vanish.

If $(L,m)=(4,2)$ then $u=2$ and $s=1$, and \eqref{eq:target} reads
$[z^{2j+1}]Q=f\binom{1}{j}$. For $j=0$ this is $[z^{1}]Q=f$ and for $j=1$ it is
$[z^{3}]Q=[z^{L-1}]Q=f$, both by Lemma~\ref{lem:Q}; for $j\ge2$ we have $2j+1>4=\deg Q$ and
both sides vanish.
\end{proof}

\begin{remark}
The argument is short because Theorem~\ref{thm:structure} has converted the criterion into
two properties of the single polynomial $Q$: it is palindromic of degree $L$ with linear
coefficient $f$, and the residues attainable subject to $\gcd(n,k)=1$ are exactly
$\{1,L-1\}$. The second is the crystallographic restriction. At no point does the proof need
to know which of $1,2,3,4,6$ the integer $L$ is.
\end{remark}

\section{The crystallographic restriction}\label{sec:crystal}

The plane crystallographic restriction says that a lattice in $\R^2$ admits rotational
symmetry only of order $1,2,3,4$ or $6$, because the trace $2\cos(2\pi/f)$ of the rotation
must be a rational integer, i.e.\ $\varphi(f)\le2$. We now show that the same five numbers
are forced here, by a different rationality constraint.

\begin{lemma}[which residues occur]\label{lem:residues}
Let $g=\gcd(m,n)$ and $0\le s<m$. There exists $k$ with $k\equiv s\pmod m$ and
$\gcd(k,n)=1$ if and only if $\gcd(s,g)=1$.
\end{lemma}

\begin{proof}
If a prime $p$ divides both $s$ and $g$ then $p\mid n$ and $p\mid k$ for every $k\equiv
s\pmod m$, so no such $k$ is coprime to $n$. Conversely assume $\gcd(s,g)=1$ and let $p\mid
n$ be prime. If $p\mid m$ then $p\mid g$, so $p\nmid s$ and hence $p\nmid k$ automatically
for every $k\equiv s\pmod m$. If $p\nmid m$ then the progression $s+m\Z$ meets every residue
class modulo $p$, so we may avoid $0$; the finitely many such conditions are simultaneously
satisfiable by the Chinese remainder theorem.
\end{proof}

\begin{lemma}[uniformity constrains primes]\label{lem:uniformprime}
Let $v$ be uniform of order $m$, and let $p$ be a prime with $p\mid n$ and $p\nmid f$. Then
$\ord_p(v)=m$; in particular $m\mid p-1$ and $m\le p-1$.
\end{lemma}

\begin{proof}
Let $H=(n/p)\Zn\cong\Z/p\Z$, a $\langle v\rangle$-stable subgroup, and take $x=(n/p)u\in H$
with $p\nmid u$. Then $v^{j}x=x$ iff $n\mid (v^{j}-1)(n/p)u$ iff $p\mid v^{j}-1$, so the
$\langle v\rangle$-orbit of $x$ has size $\ord_p(v)$. If $x$ were fixed then $p\mid v-1$ and
$p\mid n$, so $p\mid\gcd(v-1,n)=f$, contrary to hypothesis. Hence $x$ is not fixed, and
uniformity forces its orbit size to be $m$, i.e.\ $\ord_p(v)=m$. As $\ord_p(v)$ divides
$p-1$, the claim follows.
\end{proof}

\begin{lemma}[arithmetic]\label{lem:arith}
Let $f\ge1$ be an integer such that every prime $q\le f$ with $q\ne f-1$ divides $f$. Then
$f\in\{1,2,3,4,6\}$.
\end{lemma}

\begin{proof}
The values $f=1,2,3,4,6$ satisfy the condition: for $f=6$ the primes $\le6$ are $2,3,5$ and
$5=f-1$ is exempt, while $2\mid6$ and $3\mid6$; the others are immediate. For $f=5$ the
primes $\le5$ are $2,3,5$ and $f-1=4$ is not prime, so $2\mid5$ would be required, which
fails.

Now suppose $f\ge7$. Every prime $q\le f/2$ satisfies $q\ne f-1$, because $f-1>f/2$ for
$f>2$; hence every prime $q\le f/2$ divides $f$, and therefore
\[
  f\;\ge\;\prod_{p\le f/2}p\;=\;e^{\theta(f/2)},
\]
where $\theta$ is Chebyshev's function. Iterating from $f\ge7$: primes $2,3\le f/2$ give
$6\mid f$, so $f\ge12$; then $2,3,5\le f/2$ give $30\mid f$, so $f\ge30$; then
$2,3,5,7,11,13\le f/2$ give $30030\mid f$, so $f\ge30030$. For $f\ge30030$ the Chebyshev
bound $\theta(x)>0.84x$ (valid for $x\ge101$) yields $f\ge e^{0.42f}$, which is false. Hence
no $f\ge7$ satisfies the condition.
\end{proof}

\begin{theorem}[necessity of the crystallographic restriction]\label{thm:necessity}
Let $v$ be a uniform multiplier of order $m$ with $f<m$, and assume
Conjecture~\ref{conj:irr}. If $v$ sieves, then
\[
  f\in\{1,2,3,4,6\}\qquad\text{and}\qquad \rad(f)\mid m .
\]
When $f=m-1$ the appeal to Conjecture~\ref{conj:irr} may be dropped.
\end{theorem}

\begin{proof}
Uniformity gives $n=f+ma$, where $a$ is the number of orbits of size $m$; since $f<m$ we
have $r:=n\bmod m=f$, and moreover $P=\Fix(v)$, $Q(z)=(1+z)^{f}$ and $n_m=a$. Let
$s=k\bmod m$ and $b=\lfloor k/m\rfloor$; then $b\le\lfloor(n-1)/m\rfloor=a$, so
$\binom{a}{b}\ge1$. Exactly as in Case (a) of the proof of Theorem~\ref{thm:sufficiency} ---
an argument which uses only $\deg Q=f<m$, and not the criterion --- identities
\eqref{eq:expand} and \eqref{eq:eval} give
\[
  [z^k]N=\binom{a}{b}\binom{f}{s},
  \qquad
  f\,Y_{n,k}(\zeta_m)=f\binom{a}{b}\rho_{f,s},
\]
so by \eqref{eq:split} the sieving identity at $k$ is equivalent to
\begin{equation}\label{eq:necessary}
  \rho_{f,s}\;=\;\binom{f}{s}\Big/f .
\end{equation}
In particular $\rho_{f,s}$ must be \emph{rational} for every attainable $s$. We show that no
$s$ in the bad set $B=B(f)=\{0,2,3,\dots,f-2,f\}$ is attainable; note that $B$ is empty
unless $f\ge2$.

For $s\in\{0,f\}$ we have $\rho_{f,s}=1/[f]_\zeta$, so \eqref{eq:necessary} would force
$[f]_\zeta=f\in\Q$, contradicting Corollary~\ref{cor:qintirr} since $2\le f<m$. This case
is unconditional.

For $2\le s\le f-2$ --- a range non-empty only when $f\ge4$ --- the right-hand side of
\eqref{eq:necessary} obeys
\[
  \binom{f}{s}\Big/f\;\ge\;\binom{f}{2}\Big/f\;=\;\frac{f-1}{2}\;>\;1 .
\]
If $f\le m-2$, then $\rho_{f,s}$ is irrational by Conjecture~\ref{conj:irr} and
\eqref{eq:necessary} fails. If $f=m-1$, then $|\rho_{f,s}|=1$ by
Proposition~\ref{prop:rlast}, and \eqref{eq:necessary} fails on absolute values. In either
case $s$ is not attainable.

Hence no $k$ coprime to $n$ has $k\bmod m\in B$. By Lemma~\ref{lem:residues} this says
$\gcd(s,g)>1$ for every $s\in B$, where $g=\gcd(m,n)$.

In particular, for every prime $q$ with $2\le q\le r=f$ and $q\ne f-1$ we have $q\in B$ and
therefore $q\mid g$; so $q\mid m$ and $q\mid n$. Suppose such a $q$ did not divide $f$. Then
Lemma~\ref{lem:uniformprime} applies to the prime $q$ and gives $m\le q-1$; but $q\mid m$
gives $q\le m$, so $q\le m\le q-1$, a contradiction. Hence every prime $q\le f$ with
$q\ne f-1$ divides $f$, and Lemma~\ref{lem:arith} yields $f\in\{1,2,3,4,6\}$.

Finally we check $\rad(f)\mid m$. Let $q$ be any prime dividing $f$. Then $q\ne f-1$:
otherwise $q$ would divide both $f$ and $f-1$, hence $q\mid1$. So $q$ falls under the case
already treated, giving $q\mid g\mid m$. As this holds for every prime divisor of $f$, we
conclude $\rad(f)\mid m$.
\end{proof}

\begin{remark}[where the false hypothesis was felt]\label{rem:99}
The step that would fail without Proposition~\ref{prop:rlast} is the treatment of
$2\le s\le f-2$ when $f=m-1$. The smallest multiplier at which it is actually needed is
\[
  n=99,\qquad v=19,
\]
for which $m=10$, $f=9=m-1$, $v$ is uniform, and $s=3$ is attainable --- take $k=13$. The
earlier argument would have asserted $Y_{99,13}(\zeta_{10})\notin\Q$; in truth
$\rho_{9,3}=-1$ by Proposition~\ref{prop:rlast}, so that
\[
  Y_{99,13}(\zeta_{10})\;=\;\binom{9}{1}\rho_{9,3}\;=\;-9,
\]
a rational --- indeed negative --- integer, whereas the true number of fixed necklaces is
$\binom91\binom93/9=84$. Sieving does fail here, as it must, but on numerical rather than
rationality grounds. We note that $99>96$, the largest $n$ reached by the verification of
Section~\ref{sec:verification} at the time the hypothesis was formulated; neither the proof
nor the computation could have detected the gap.
\end{remark}

\begin{remark}
Theorem~\ref{thm:necessity} covers uniform multipliers with $f<m$; together with
Theorem~\ref{thm:sufficiency} it settles that case completely, modulo
Conjecture~\ref{conj:irr}. When $f\ge m$ --- in particular when $m\mid n$, where
Proposition~\ref{prop:lucas} degenerates and Proposition~\ref{prop:mdivn} must be used
instead --- the same conclusion is observed in every computed case, but the argument above
does not apply verbatim. Conjecture~\ref{conj:converse} for non-uniform multipliers remains
open.
\end{remark}

\begin{remark}[why the coincidence is not one]
Both the lattice statement and Theorem~\ref{thm:necessity} amount to demanding that a
quantity built from a root of unity be rational: for a lattice it is the trace of a
rotation, here it is a Gaussian binomial divided by a $q$-integer. In each case the demand
confines the root of unity to a field of degree at most two over $\Q$, and the surviving
orders are $\{1,2,3,4,6\}$. A necklace contains no lattice; what the two situations share is
only the circle.
\end{remark}

\section{Computational verification}\label{sec:verification}

All computations use exact integer arithmetic, together with exact arithmetic in
$\Z[q]/(\Phi_m(q))$ for root-of-unity evaluations. No floating point is used anywhere.

\paragraph{Method.} The polynomial side is computed from Proposition~\ref{prop:factor} as a
product of cyclotomic polynomials reduced modulo $\Phi_m$; the fixed-point side from
Corollary~\ref{cor:twocolour}. Both were cross-validated against independent routes: the
factorisation against direct polynomial division of $\qbin{n}{k}$ by $\qint{n}$ for all
$n<30$, and the fixed-point formula both against the affine sum of Lemma~\ref{lem:disjoint}
over all $c\in\Zn$ and against brute-force enumeration of necklaces. Rationality of an
element of $\Z[q]/(\Phi_m)$ is decided by comparing it with a rational multiple of the
reduced denominator, so no root of unity is ever approximated.

\paragraph{Scope.} Each multiplier is tested against \emph{every} admissible $k$; the
``instances'' column counts multipliers for the criterion rows and individual checks
otherwise.

\begin{table}[h]
\centering
\small
\begin{tabular}{llrr}
Statement & range & instances & counterex.\\\hline
Thm.~\ref{thm:collapse}, two colours, criterion \emph{not} assumed & $n\le 60$ & $29{,}780$ & $0$\\
Thm.~\ref{thm:collapse}, three colours, every content & $n\le 28$ & $34{,}186$ & $0$\\
Thm.~\ref{thm:structure}, structure of $P$ & $n\le 400$ & $19{,}820$ & $0$\\
Coefficient identities of \S\ref{sec:sufficiency} & $n\le 200$ & $1{,}048{,}609$ & $0$\\
Sieving identity, criterion satisfied (Thm.~\ref{thm:sufficiency}) & $n\le 130$ & $2{,}527$ & $0$\\
Sieving identity, criterion violated (Conj.~\ref{conj:converse}) & $n\le 130$ & $2{,}497$ & $0$\\
Uniform $(m,f)$ form (Cor.~\ref{cor:uniform}) & $n\le 115$ & $2{,}682$ & $0$\\
Cor.~\ref{cor:qintirr}, $[r]_\zeta\notin\Q$ & $m\le 60$ & $1{,}711$ & $0$\\
Prop.~\ref{prop:conj}, congruence necessary & $m\le 60$ & $34{,}278$ & $0$\\
Conj.~\ref{conj:irr}, restricted irrationality & $m\le 60$ & $34{,}278$ & $0$\\
\end{tabular}
\caption{Exhaustive verification. The rows for Theorems~\ref{thm:collapse}
and~\ref{thm:structure} test statements proved above and so are consistency checks on the
implementation; the row for Conjecture~\ref{conj:converse} is genuine evidence, as is the
last row. Every counterexample to the discarded irrationality hypothesis found in the search
of Remark~\ref{rem:hypfalse} has $r=m-1$, and there are $67$ of them with $m\le60$.}
\label{tab:verification}
\end{table}

\section{Boundaries and open problems}\label{sec:open}

\paragraph{Uniformity is not necessary.} Of $280$ non-uniform multipliers with $n\le60$,
exactly $38$ satisfy the sieving identity; all of them have orbit-size profile
$\{1,2,m\}$. Theorem~\ref{thm:structure} explains that profile: a non-uniform multiplier
meeting the criterion must act on $C_L$ as $-1$, and the orbits of $-1$ on a cyclic group
have size $1$ or $2$. This is also why Definition~\ref{def:criterion} ranges over all
divisors $e\mid m$ rather than taking the $(m,f)$ form: the single-level condition on
$f=f_1$ is necessary but not sufficient, failing in $106$ of $603$ multipliers with
$n\le46$.

\paragraph{Three colours is strictly stronger.} For contents with $\ell\ge3$ colours the
criterion as stated is no longer correct: among $213$ multipliers with $n\le28$ there are
$25$ for which the two-colour identity holds but the three-colour identity fails --- the
smallest at $n=28$. In every discrepancy the criterion over-predicts. Theorem~\ref{thm:collapse}
holds verbatim for every content, so the counting side of the identity is unaffected and the
whole discrepancy lies on the $q$-analogue side: it is the choice of $q$-multinomial, not the
enumeration of fixed necklaces, that fails to generalise. The additional obstruction is not
understood.

\paragraph{Open problems.}
\begin{enumerate}
  \item Prove Conjecture~\ref{conj:irr}, or the weaker modulus statement of
        Remark~\ref{rem:modulusone}, which would render Theorem~\ref{thm:necessity}
        unconditional.
  \item Prove Conjecture~\ref{conj:converse} for non-uniform multipliers, and extend
        Theorem~\ref{thm:necessity} to $f\ge m$.
  \item Determine the correct criterion for contents with $\ell\ge3$ colours, given that by
        Theorem~\ref{thm:collapse} the obstruction is purely $q$-analytic.
\end{enumerate}

\appendix

\section{An unrelated side result: orbit towers terminate}\label{app:tower}

A sequence $a=(a_j)_{j\ge1}$ of non-negative integers is \emph{realisable} if
$L(a)_j=\frac1j\sum_{d\mid j}\mu(j/d)\,a_d$ is again a sequence of non-negative integers;
then $a_j$ counts the points of period $j$ of a dynamical system and $L(a)_j$ counts its
orbits of length $j$ --- its ``primes''. One may ask whether the primes of a system can
themselves be the points of a system, indefinitely.

\begin{theorem}\label{thm:tower}
Let $a\ne0$ be a sequence of non-negative integers. Then $L^{\,j}(a)$ fails to be a sequence
of non-negative integers for some $j\ge1$.
\end{theorem}

\begin{proof}
Let $N=\min\{j: a_j\ne0\}$. If $N\ge2$ then $a_j=0$ for $j<N$, and this persists: if
$c_j=0$ for all $j<N$ then $L(c)_j=0$ for $j<N$, while $L(c)_N=c_N/N$ since every proper
divisor of $N$ contributes $0$. Hence $L^{\,j}(a)_N=a_N/N^{\,j}$, which is a positive
integer for all $j$ only if $a_N=0$, a contradiction.

If $N=1$, put $\lambda=a_1>0$ and $x_j=L^{\,j}(a)_2$. Since $L$ fixes the first term,
$L^{\,j}(a)_1=\lambda$ for all $j$, and $x_{j+1}=(x_j-\lambda)/2$. Solving,
$x_j=-\lambda+(a_2+\lambda)2^{-j}\to-\lambda<0$, so $x_j<0$ for large $j$.
\end{proof}

Thus the tower of ``primes of the primes of the primes'' is always finite, with height at
most $\log_N(a_N)$ when $N\ge2$. Note the asymmetry: the opposite direction,
$a\mapsto(\sum_{d\mid j}d\,a_d)_j$, may be iterated forever.

\section{Attribution ledger}\label{app:ledger}

The results above emerged from a broad computational search, most of whose output was
already known. Recording that is part of the result.

\begin{itemize}
  \item \textbf{New.} The criterion of Definition~\ref{def:criterion}, its sufficiency
        (Theorem~\ref{thm:sufficiency}) and Theorem~\ref{thm:necessity}; the cases
        $f\in\{3,4,6\}$ are not covered by \cite{stucky}, and lie exactly in the regime that
        paper flags as having no general conjecture.
  \item \textbf{New.} Theorem~\ref{thm:collapse}, in particular the vanishing of
        $N_\alpha(A_c)$ for $c\notin(v-1)\Zn$ and the resulting $f_1$-to-one correspondence,
        valid for every colour content.
  \item \textbf{New.} Theorem~\ref{thm:structure}, which recasts the criterion as a
        statement about a single subgroup of crystallographic order.
  \item \textbf{Corrected here.} The irrationality hypothesis assumed in the first version of
        this work is false; see Remark~\ref{rem:hypfalse}. Proposition~\ref{prop:conj},
        Corollary~\ref{cor:qintirr} and Proposition~\ref{prop:rlast} replace it, and
        Remark~\ref{rem:99} records the first multiplier at which the difference is felt.
  \item \textbf{Rediscovered.} Sieving for prime $n$: derived and proved here via $q$-Lucas
        before being located as \cite[Cor.~3.3]{stucky}. The result is Stucky's.
  \item \textbf{Rediscovered.} That
        $(\qint{n}/\qint{n-mk})\qbin{n-mk}{k}$ cyclically sieves the $k$-subsets of a cycle
        with all cyclic gaps $\ge m+1$ (verified to $n=16$).
  \item \textbf{Rediscovered.} That the central binomial, Ap\'ery, Franel and Domb numbers
        are periodic-point counts; their orbit counts are catalogued in the OEIS as the
        \texttt{A060165} family.
  \item \textbf{Observed, uncatalogued.} Necklaces equal to their own Burrows--Wheeler
        transform. For $n=2,\dots,18$ their numbers are
        \[3,4,4,4,7,6,5,9,5,6,9,8,8,11,11,7,15,\]
        so among the $4116$ necklaces of length $16$ exactly $11$ are BWT-fixed. We found no
        database entry for this sequence and offer no characterisation.
\end{itemize}

\begin{thebibliography}{9}
\bibitem{rsw} V.~Reiner, D.~Stanton and D.~White,
  \emph{The cyclic sieving phenomenon},
  J.\ Combin.\ Theory Ser.\ A \textbf{108} (2004), 17--50.
\bibitem{stucky} E.~N.~Stucky,
  \emph{Parity-Unimodality and a Cyclic Sieving Phenomenon for Necklaces},
  arXiv:1812.04578.
\bibitem{sagan} B.~E.~Sagan,
  \emph{The cyclic sieving phenomenon: a survey},
  in: Surveys in Combinatorics 2011, London Math.\ Soc.\ Lecture Note Ser.\ \textbf{392},
  Cambridge Univ.\ Press, 183--233.
\end{thebibliography}

\end{document}
